\begin{textsample}{POS Dim 4 – persona\_grok – Score 17.00 – t558\_grok.txt}  \label{ex:f4_pos_034}
Greenpeace 's Pole to Pole voyage took us on an incredible journey across the world 's \textbf{oceans}, where we \textbf{documented} breathtaking wonders alongside urgent \textbf{threats} from climate change, overfishing, plastic pollution, and deep-sea \textbf{mining}.
Throughout this expedition, we advocated strongly for a UN Global Ocean Treaty that would establish \textbf{sanctuaries} to protect these vital \textbf{ecosystems}.
Our travels began in the Antarctic, where we \textbf{witnessed} the impacts on \textbf{penguin} populations, with \textbf{declines} linked to scarcity of krill, a key food source affected by environmental pressures.
Moving to the Amazon Reef, our team confirmed the presence of \textbf{whale} breeding grounds, highlighting the \textbf{area} 's importance for marine \textbf{life}.
In the Sargasso Sea, we observed \textbf{wildlife} entangled in plastic debris, underscoring the pervasive pollution crisis.
Finally, at the Lost City hydrothermal vents, we explored unique underwater formations now at risk from potential deep-sea \textbf{mining} \textbf{activities}.
As part of our efforts, we produced the film"Turtle Journey"to raise awareness about \textbf{ocean} conservation.
We also directly confronted \textbf{fishing} \textbf{vessels} operating in sensitive \textbf{areas}, and our researchers gathered data on \textbf{turtles} and \textbf{sharks} to support \textbf{protection} initiatives.
Despite these discoveries and actions, UN negotiations on \textbf{ocean} \textbf{protection} have been disappointing so far.
That 's why we 're urging leaders and communities worldwide to push for stronger global action to \textbf{safeguard} our \textbf{oceans} for future generations.

% matched lemmas: activity, area, decline, document, ecosystem, fishing, life, mining, ocean, penguin, protection, safeguard, sanctuary, shark, threat, turtle, vessel, whale, wildlife, witness
\end{textsample}
